% !TEX root = ../main.tex
% Export dei diagrammi in docs/diagrams/use-cases/exported/.
\subsection{Diagrammi dei casi d'uso per attore}
\label{sec:diagrammi-attori}
\begingroup
\setlength{\abovecaptionskip}{4pt}
\setlength{\belowcaptionskip}{0pt}
\begin{figure}[!htbp]
  \centering
  \includegraphics[width=0.32\textwidth]{use-cases/use-cases-visitor.png}
  \caption{Casi d'uso associati al Visitor.}
  \label{fig:uc-visitor}
\end{figure}
\begin{figure}[!htbp]
  \centering
  \includegraphics[width=0.88\textwidth]{use-cases/use-cases-user.png}
  \caption{Casi d'uso associati allo User.}
  \label{fig:uc-user}
\end{figure}
\begin{figure}[!htbp]
  \centering
  \includegraphics[width=0.94\textwidth]{use-cases/use-cases-organiser-administrator.png}
  \caption{Casi d'uso associati a Organiser e Administrator.}
  \label{fig:uc-organiser-administrator}
\end{figure}
\FloatBarrier
\endgroup
